\documentclass[12pt]{article}
\usepackage[T1]{fontenc}
\usepackage[utf8]{inputenc}
\usepackage{lmodern}
\usepackage{textcomp}
\usepackage{enumitem}
\usepackage{geometry}
\geometry{margin=1in}
\begin{document}
\begin{center}
Answer Key - Chemistry - Graduate Level Assessment 8 \
\small (Answers are ordered exactly as in the question paper.)
\end{center}

\section*{Multiple Choice}
\begin{enumerate}[label=\arabic*.]
\item \textbf{Answer:} A
\\ \textit{Options:} A) 22$\mathrm{atm}$; B) 20$\mathrm{atm}$; C) 15$\mathrm{atm}$; D) 30$\mathrm{atm}$; E) 34$\mathrm{atm}$
\\ \textit{Note:} The correct answer of 22 atm arises from applying the van der Waals equation, which accounts for the non-ideal behavior of gases by including factors for intermolecular forces and molecular volume, thus providing a more accurate pressure calculation for xenon gas under the specified conditions.
\item \textbf{Answer:} B
\\ \textit{Options:} A) 0.9; B) 1.0; C) 2.5; D) 0.3; E) 0.5
\\ \textit{Note:} The equilibrium constant (K) for the reaction H$_{2}$ + I$_{2}$ ⇌ 2 HI is calculated using the formula K = [HI]^2 / ([H$_{2}$][I$_{2}$]), and substituting the given equilibrium concentrations yields K = (0.$6)^2$ / (0.9 * 0.4) = 1.0.
\item \textbf{Answer:} A
\\ \textit{Options:} A) 2 $10^{-10} \mathrm{~m}$; B) 1 $10^{-10} \mathrm{~m}$; C) 7 $10^{-10} \mathrm{~m}$; D) 6 $10^{-10} \mathrm{~m}$; E) 4 $10^{-10} \mathrm{~m}$
\\ \textit{Note:} The correct answer of \(2 \times 10^{-10} \mathrm{~m}\) for \(\lambda_{\max}\) corresponds to the peak wavelength of radiation emitted by a black body at a temperature of \(10^7 \mathrm{~K}\), which can be calculated using Wien's displacement law, \(\lambda_{\max} = \frac{b}{T}\) where \(b\) is Wien's displacement constant (\(2.898 \times 10^{-3} \mathrm{~m \cdot K}\)).
\item \textbf{Answer:} A
\\ \textit{Options:} A) 169 $\mathrm{g} \mathrm{mol}^{-1}$; B) 123 $\mathrm{g} \mathrm{mol}^{-1}$; C) 200 $\mathrm{g} \mathrm{mol}^{-1}$; D) 150 $\mathrm{g} \mathrm{mol}^{-1}$; E) 190 $\mathrm{g} \mathrm{mol}^{-1}$
\\ \textit{Note:} The correct answer of 169 g mol⁻¹ is derived from the ideal gas law, using the formula M = (density × R × T) / P, where R is the ideal gas constant, T is the temperature in Kelvin, and P is the pressure in pascals, leading to the calculation of molar mass based on the given density, temperature, and pressure conditions.
\item \textbf{Answer:} C
\\ \textit{Options:} A) lead; B) cesium; C) bromine; D) francium; E) rubidium
\\ \textit{Note:} Bromine is the only liquid element that is dark-colored and nonconducting at room temperature, distinguishing it from other elements and fitting the criteria described in the question.
\end{enumerate}

\section*{True / False}
\begin{enumerate}[label=\arabic*.]
\setcounter{enumi}{5}
\item \textbf{Answer:} True
\\ \textit{Options:} A) True; B) False
\\ \textit{Note:} The statement is true because 600 mm of water is equivalent to approximately 44.1 mm of mercury due to the difference in density between water and mercury, with water being around 13.6 times less dense than mercury.
\item \textbf{Answer:} True
\\ \textit{Options:} A) True; B) False
\\ \textit{Note:} The statement is true because the mass of a proton is approximately 1.67 x 10^-27 kilograms, while the mass of a positron is about 9.11 x 10^-31 kilograms, making the proton significantly more massive than the positron.
\item \textbf{Answer:} False
\\ \textit{Options:} A) True; B) False
\\ \textit{Note:} The statement is false because the logarithm of a positive number, such as 0.0364, cannot be negative and is instead a negative value, specifically the logarithm is approximately -1.4380, not -1.5000.
\item \textbf{Answer:} False
\\ \textit{Options:} A) True; B) False
\\ \textit{Note:} The original energy level of the electron for transitions in the Lyman series must be 2 or higher, as the Lyman series corresponds to transitions from higher energy levels (n ≥ 2) to the lowest energy level (n = 1), making the statement false.
\item \textbf{Answer:} True
\\ \textit{Options:} A) True; B) False
\\ \textit{Note:} The statement is true because the complete combustion of butane (C$_{4}$H$_{10}$) requires a stoichiometric ratio of 1 mole of butane to 13 moles of oxygen, which, when calculated, shows that 1000 cubic feet of butane indeed requires 6500 cubic feet of oxygen for complete combustion.
\end{enumerate}

\section*{Long Form Response}
\begin{enumerate}[label=\arabic*.]
\setcounter{enumi}{10}
\item \textbf{Answer:} To calculate the volume of concentrated sulfuric acid required to achieve a mass of 74.0 g, we can use the relationship between mass, volume, and density, which is given by the formula: mass = density \ensuremath{\times} volume. Rearranging this formula to solve for volume, we find volume = mass / density. Substituting the given values, we have volume = 74.0 g / 1.85 g/ml, which equals approximately 40.0 ml. Therefore, to achieve a mass of 74.0 g, we need about 40.0 ml of concentrated sulfuric acid. This calculation demonstrates the direct application of density in determining the necessary volume based on the desired mass of a substance.
\\ \textit{Note:} correct because it accurately applies the relationship between mass, volume, and density, using the formula volume = mass / density to determine that approximately 40.0 ml of concentrated sulfuric acid is needed to achieve a mass of 74.0 g, given its density of 1.85 g/ml.
\item \textbf{Answer:} The density of the object can be calculated using the formula density = mass/volume. Given the mass of the object is 200 g and the volume of water displaced is 60 ml, the density of the object is 200 g / 60 ml = 3.33 g/ml. For the oil, which displaced 50 g of volume, we can determine its density by considering that 1 ml of water is approximately equivalent to 1 g; thus, the displaced volume of the oil is 50 ml. Consequently, the density of the oil is 50 g / 60 ml = 0.83 g/ml. The significance of these findings lies in the comparison of densities: the object is denser than both water and oil, indicating that it will sink in both mediums, while the oil, being less dense than water, will float, illustrating the principles of buoyancy and density in fluid mechanics.
\\ \textit{Note:} The answer correctly applies the principle of displacement to calculate the densities of both the object and the oil by using the mass-to-volume ratio, revealing that the object is denser than water, while the oil is less dense, which is significant in understanding buoyancy and material properties.
\end{enumerate}

\vfill
\noindent\textit{End of Answer Key}
\end{document}